\graphicspath{{../02Physics/pics/}}
	
\chapter[Physics]{Physics}\label{ch:Physics}

\lettrine[lines=2]{\color{darkocre}N}{umbers} are powerful
mathematical objects. They are used to solve
an endless list of problems that involve \emph{quantities}. As
mathematics
and sciences progressed, natural numbers evolved into whole
numbers, then into rational numbers and beyond.\footnote{A superb account of
this process is given in the book \emph{``Number: The Language of
Science''} by Tobias Dantzig.}

\begin{myprereq}{Prerequisite Knowledge}
	To fully understand the material of this chapter, readers should be comfortable with the following concepts:
	
	\begin{itemize}
		\item \phantom{phantom}
		\vspace{-0.5cm}
		\item State
		\item Dynamical equations
	\end{itemize}	
\end{myprereq}

\section{Goals and Methods}
Physics is a human activity pursuing the following major goals: \emph{Describe}, \emph{explain}, and \emph{predict} phenomena that comprise the observed world.

results can be applied in a wide range of fields. In part, the universality of mathematics
stems from the \emph{general} and \emph{abstract} nature of mathematical
concepts. Let us illustrate this using an example.

\begin{SCfigure}%[htbp]
  %\centering
  \includegraphics[scale=1.0]{defaultFigureTemplate}
  \caption{49 objects can be arranged in a square 7x7. 48 objects can
    be arranged as a rectangle of 6x8.}
  \label{fig:numbersExampleGenerality}
\end{SCfigure}

An astute farmer notices that 49 sacks of grains can be arranged
in a square with each side having 7 sacks (see the Figure
\ref{fig:numbersExampleGenerality}). When one sack is used up, the
remaining 48 sacks can be arranged as a rectangle 6 by 8 sacks.


\begin{exercise}\label{exe:relationsGeneral}
Think how you would represent the generalized relations of the types
given in the Figure \ref{fig:schematicRelationNtoN} at the level of
sets? What kind of diagrams would you draw?
\end{exercise}

\section{Common Sense}
Mathematics is a remarkably effective and universal discipline, its
methods and
results can be applied in a wide range of fields.

\section{Determinism}
Mathematics is a remarkably effective and universal discipline, its
methods and
results can be applied in a wide range of fields.

\section{Classical and Quantum}
Mathematics is a remarkably effective and universal discipline, its
methods and
results can be applied in a wide range of fields.

\section{Atoms}
Classical physics predicts a continuous decay of unstable configuration of charges. What is observed is a spontaneous decay of stable configuration of charges. Quantum physics elegantly explains the latter.  

\section{Particles}
Classical physics predicts a continuous decay of unstable configuration of charges. What is observed is a spontaneous decay of stable configuration of charges. Quantum physics elegantly explains the latter.  

\section{Polarization and Spin}
Mathematics is a remarkably effective and universal discipline, its
methods and
results can be applied in a wide range of fields.

% ---------------------- Chapter Highlights -------------------------------
\section*{Chapter Highlights}
{\setstretch{1.5}\chhc
  \it  
\begin{itemize}
\item The power of mathematical concepts and methods increases with
  the level of abstraction.
\item Learning new concepts often involves learning new
  terminology. The latter can create an artificial mental barrier.
\item ``Usual'' numbers form a mathematical structure. The structure
  is revealed through various relations that exist between numbers.
\item Relations between numbers are expressed using the concept of
  functions and operations (e.g., addition). Each operation is
  characterized by its arity -- the number of arguments it accepts as
  an input.
\end{itemize}


}
